\section{Methodology}\label{sec:method}

We separate two questions that are easy to conflate: how the model represents the \emph{drama construct} itself, and how it represents \emph{reader variation} in drama. Our experiments are designed so that one set isolates the geometry of the construct (E1, E2, E4) and a second set isolates the geometry of disagreement (E3). We describe the hypotheses, the model and extraction pipeline, the two data sources, the five experiments, and the metrics and controls.

\para{Hypotheses.} We test two accounts of how \qwenbig encodes what different people find dramatic. Under \textbf{H1}, drama is a low-dimensional shared feature space, and readers differ by \emph{thresholds or regions} on that space rather than by what they read; in the single-axis limit this reduces to one drama axis plus a per-reader threshold (intercept). Under \textbf{H2}, readers differ by a \emph{reader-specific direction} inside a multi-dimensional drama subspace, so that two people can disagree because they project the same activations onto different drama features. These hypotheses concern two distinct geometries, and we keep them apart on purpose: the construct can be multi-dimensional while reader variation is still a one-dimensional threshold effect, which is in fact what we find. We pre-registered the decision rule. \textbf{H2} is supported \emph{iff} (i) the melodrama residual is decodable above a control \emph{and} (ii) the held-out AUROC gain from giving each reader their own direction, $\Delta\mathrm{AUROC}(\Mthree-\Mtwo)$, has a bootstrap 95\% CI that excludes zero. \textbf{H1} is supported \emph{iff} \Mone${\approx}$\Mtwo${\approx}$\Mthree and the melodrama residual matches its control. The construct-geometry sub-hypotheses (single axis vs.\ multi-dimensional space) are evaluated separately from the reader-variation sub-hypotheses (threshold vs.\ direction).

\para{Models and activation extraction.} Our primary model is \qwenbig (28 layers, hidden size 3584); we replicate the full pipeline on \qwensmall (28 layers, hidden size 1536) as a scale ablation. We run both from HuggingFace in \texttt{bfloat16} on a single RTX A6000; \texttt{fp16} produced NaNs in the deep layers, which we document and avoid. For each text we take hidden states and mean-pool over the text tokens, extracting at layers spanning relative depths $\{0.25, 0.5, 0.57, 0.75, 1.0\}$. We use raw text for all base-geometry experiments (E1--E3) and the chat template with a system persona for the persona probe (E4). All concept directions are difference-in-means vectors, which are robust and causally meaningful~\citep{marks2023geometry}: for a binary contrast with positive and negative example sets, the direction is
\[
\mathbf{v} \;=\; \boldsymbol{\mu}_{+} - \boldsymbol{\mu}_{-},
\qquad
\boldsymbol{\mu}_{\pm} = \tfrac{1}{|S_{\pm}|}\sum_{x \in S_{\pm}} h_\ell(x),
\]
where $h_\ell(x)$ is the mean-pooled activation at layer $\ell$. We fix all seeds to 42. APIs are used only for stimulus generation and validation; every model-internal computation runs locally on the GPU.

\para{Controlled drama corpus (\dramacorpus).} To probe construct geometry under tight control, we build \dramacorpus. It contains 24 base scenarios crossed with 4 registers $\{$neutral, understated, dramatic, melodramatic$\}$ for 96 register-matched texts, plus 8 drama subconcepts (suspense, pathos/grief, spectacle, stakes, sincerity, excess/sentimentality, shock, restraint), each with ${\sim}24$ contrast texts, for 194 additional texts. The corpus is produced by a 28-agent generation workflow and then adversarially blind re-labeled by independent agents; we keep only label-agreed items. All 290 of 290 items passed at 100\% agreement, and the kept set is perfectly register-balanced, so the dramatic-vs-melodramatic contrast is not confounded by scenario or length.

\para{Real reader labels (\goemotions).} For real disagreement we use raw \goemotions: 211k annotations from 70 identified raters. We keep texts with $\geq 3$ raters and raters with $\geq 150$ labels, yielding 56,995 texts. We operationalize drama as perceived emotional intensity/arousal: for each (text, rater) pair, $\mathrm{drama\_label}=1$ \emph{iff} the rater chose any high-arousal emotion (arousal $\geq 0.70$ under the NRC-VAD circumplex mapping). The resulting base rate is $0.185$, and $36\%$ of texts are \emph{contested} (rater drama-rate in $[0.15, 0.85]$), confirming substantial real disagreement to model rather than label noise. We inspect \emobank as an aggregate-only auxiliary; its individual ratings are anonymized, so it is not central to the reader-variation analysis.

\para{Experiments.} \textbf{E1} fits the drama axis on \dramacorpus and asks {\bf is melodrama just more drama?} by measuring the cosine and principal angle between the drama axis and the dramatic$\rightarrow$melodramatic step across depths. \textbf{E2} forms the difference-in-means directions for the 8 subconcepts and computes the participation ratio (PR) of the resulting subspace, benchmarked against PR${\approx}8$ for random directions and PR${=}1$ for a single shared axis. \textbf{E3} is the crux: a ladder of nested per-reader logistic models predicting $\mathrm{drama\_label}$ from text activations, evaluated on text-disjoint held-out data across all 70 readers. For reader $r$ with intercept $b_r$, gain $g_r$, and weight vector $\mathbf{w}_r$ over a $K$-dimensional subspace $P$, the model is
\[
\Pr(\mathrm{drama}=1 \mid h) \;=\; \sigma\!\big(g_r\,\mathbf{w}_r^{\top} P\,h \,+\, b_r\big),
\]
and the ladder progressively frees these terms: \Mzero (shared single axis, one global $\mathbf{w}$ and $b$), \Mone ($+$ per-reader threshold $b_r$), \Mtwo ($+$ per-reader gain $g_r$ on one axis), \MtwoK (shared $K$-dim direction $+$ per-reader threshold), and \Mthree (per-reader direction $\mathbf{w}_r$ in the $K$-dim subspace). We run E3 both on a \emph{transferred} \dramacorpus subspace and on an \emph{in-domain} PCA subspace ($K \in \{10, 30\}$); the in-domain version is the fair test because the subspace is fit on \goemotions activations. \textbf{E4} re-derives the drama axis under six system personas (cynic, romantic, stoic, thrill-seeker, critic, neutral) and measures the pairwise cosine between the resulting axes to ask whether a reader's stance rotates the drama direction or merely shifts where they sit on it.

\para{Metrics, baselines, and controls.} For construct geometry we report probe AUROC, participation ratio, and principal angles. For reader variation we report held-out (text-disjoint) AUROC with bootstrap 95\% CIs, and we guard the \Mthree direction effect with a \emph{shuffled-reader permutation null}: we randomize reader identity, refit \Mthree, and compare its gain over \MtwoK to the real gain. Because freeing per-reader directions adds model capacity regardless of whether reader-specific structure exists, this null is the correct baseline---any AUROC lift present under shuffled identities is pure capacity, not genuine direction. We further summarize the fitted reader-weight matrix by its rank (participation ratio over the $K$ subspace dimensions) to check whether readers actually spread across many directions or collapse onto few. Random-direction controls calibrate the melodrama-residual test in E1. Every direction, probe, and null shares the seed and extraction pipeline above, so the full ladder is reproducible end to end.
